% Artikel atau makalah singkat — templat DSWorkbench
% Kompilasi: DSWorkbench: Menulis — Kompilasi (hasil di build/main.pdf)
\documentclass[11pt,a4paper]{article}

\usepackage[utf8]{inputenc}
\usepackage[T1]{fontenc}
\usepackage[indonesian]{babel}
\usepackage[margin=2.5cm]{geometry}
\usepackage{amsmath}
\usepackage{graphicx}
\usepackage{booktabs}
\usepackage[hidelinks]{hyperref}

\graphicspath{{gambar/}}

\title{Judul Artikel yang Singkat dan Jelas}
\author{Nama Penulis\\
	\small Program Studi Sains Data, Institut Teknologi Sumatera\\
	\small \texttt{nama@contoh.ac.id}}
\date{\today}

\begin{document}
\maketitle

\begin{abstract}
	Abstrak memuat latar belakang, tujuan, metode, hasil utama, dan kesimpulan dalam satu
	paragraf, paling banyak 200 kata.

	\medskip
	\noindent\textbf{Kata kunci:} kata kunci 1, kata kunci 2, kata kunci 3
\end{abstract}

\section{Pendahuluan}
Jelaskan masalah yang dikaji, mengapa penting, dan apa yang sudah dikerjakan orang lain
\cite{penulis2024}. Tutup dengan tujuan tulisan ini.

\section{Metode}
Uraikan data dan cara kerja dengan cukup rinci agar dapat diulang orang lain. Rumus
ditulis seperti $y = \beta_0 + \beta_1 x + \varepsilon$ di dalam kalimat, atau dalam
baris tersendiri:
\begin{equation}
	\hat{\beta} = (X^{\top} X)^{-1} X^{\top} y
	\label{eq:ols}
\end{equation}

\section{Hasil dan Pembahasan}
Sajikan temuan utama, lalu bahas artinya. Tabel~\ref{tab:hasil} menunjukkan contoh
tabel.
\begin{table}[htbp]
	\centering
	\caption{Contoh tabel hasil}
	\label{tab:hasil}
	\begin{tabular}{lrr}
		\toprule
		Model     & Akurasi & Waktu (detik) \\
		\midrule
		Model A   & 0{,}91  & 12 \\
		Model B   & 0{,}88  & 3 \\
		\bottomrule
	\end{tabular}
\end{table}

\section{Kesimpulan}
Nyatakan jawaban atas tujuan tulisan dan batas-batasnya.

\begin{thebibliography}{9}
	\bibitem{penulis2024} Nama Penulis. (2024). Judul artikel. \textit{Nama Jurnal}, 1(2), 10--20.
\end{thebibliography}

\end{document}
